\documentclass[a4paper,11pt]{article}

\usepackage[margin=2.4cm]{geometry}
\usepackage{amsmath,amssymb,bm}
\usepackage[dvipsnames]{xcolor}
\usepackage{enumitem}
\usepackage{xr}
\externaldocument{PT_lectures}   % equation numbers refer to the lecture notes
\usepackage[colorlinks=true,linkcolor=RoyalBlue,citecolor=ForestGreen]{hyperref}

\newcommand{\vx}{\bm{x}}
\newcommand{\vq}{\bm{q}}
\newcommand{\vk}{\bm{k}}
\newcommand{\vu}{\bm{u}}
\newcommand{\vPsi}{\bm{\Psi}}
\newcommand{\nhat}{\hat{\bm{z}}}
\newcommand{\khat}{\hat{\bm{k}}}
\def\cH{\mathcal{H}}
\def\Plin{P_{\rm L}}
\def\knl{k_{\rm NL}}
\newcommand{\dd}{\mathrm{d}}
\def\dirac{\delta_{\rm D}}
\def\Om{\Omega_m}
\def\G{\mathcal{G}}
\def\lcdm{$\Lambda$CDM}
\def\vsv{\bm{s}}

\newenvironment{solution}[1]
  {\par\medskip\noindent\textbf{#1}\par\nobreak\smallskip}
  {\par\medskip}

\title{From Primordial Fluctuations to the Cosmic Web\\[2mm]
\large Exercise solutions}
\author{Nhat-Minh Nguyen}
\date{\today}

\begin{document}
\maketitle

\noindent Solutions to the exercises in the lecture notes. Equation numbers
refer to the notes; equations numbered here are local to each solution.

\begin{solution}{Exercise 1.1 (the Newtonian equation of motion).}
Conformal time is not an affine parameter, so we cannot use the affine form of
the geodesic equation directly. Parametrizing the worldline by $x^0 = \tau$ and
eliminating the affine parameter gives the standard result
\begin{equation*}
\ddot x^i + \Gamma^i{}_{\alpha\beta}\dot x^\alpha\dot x^\beta
- \dot x^i\,\Gamma^0{}_{\alpha\beta}\dot x^\alpha\dot x^\beta = 0 ,
\end{equation*}
with overdots denoting $\dd/\dd\tau$ and $\dot x^0 = 1$. The last term is the
one that is easy to forget; dropping it doubles the Hubble drag. For the metric
$\dd s^2 = a^2[-(1+2\Phi)\dd\tau^2 + (1-2\Phi)\dd\vx^2]$, to first order in
$\Phi$ and leading order in $u \ll 1$, the relevant symbols are
\begin{equation*}
\Gamma^i{}_{00} = \partial_i\Phi ,
\qquad
\Gamma^i{}_{0j} = \left(\cH - \dot\Phi\right)\delta^i{}_j ,
\qquad
\Gamma^0{}_{00} = \cH + \dot\Phi .
\end{equation*}
Substituting,
\begin{equation*}
\ddot x^i + 2\left(\cH-\dot\Phi\right)\dot x^i + \partial_i\Phi
- \left(\cH+\dot\Phi\right)\dot x^i = 0
\qquad\Longrightarrow\qquad
\ddot{\vx} + \cH\,\dot{\vx} = -\nabla\Phi + 3\dot\Phi\,\dot{\vx} .
\end{equation*}
The final term carries one power of the potential and one of the velocity, so it
is higher order and we drop it, leaving
\begin{equation*}
\ddot{\vx} + \cH\,\dot{\vx} = -\nabla\Phi ,
\end{equation*}
the Newtonian equation for a particle in an expanding background, with $\cH\dot{\vx}$
the Hubble drag. A quick check on the drag coefficient: with $\Phi = 0$ the
solution is $\dot{\vx} \propto 1/a$, the familiar decay of peculiar velocities.
Nothing here required $k \gg \cH$; that limit was used earlier to reduce the
Einstein constraint to the Poisson equation, which is what makes $\Phi$ the
Newtonian potential.
\end{solution}

\begin{solution}{Exercise 1.2 (the sub-horizon limit).}
\emph{(i)} The two groups in \eqref{eq:energy_constraint} are $k^2\phi$ and
$3\cH(\dot\phi + \cH\psi)$. With $\phi=\psi=\Phi$ and $\dot\Phi \lesssim \cH\Phi$,
the second is bounded by $3\cH(\cH\Phi + \cH\Phi) = 6\cH^2\Phi$, so
\begin{equation*}
\frac{\text{discarded}}{\text{retained}} \;\lesssim\; 6\left(\frac{\cH}{k}\right)^2 .
\end{equation*}
For the growing mode in matter domination the bound is saturated from below
rather than above: $\Phi$ is constant, so $\dot\Phi = 0$ exactly and the ratio is
$3(\cH/k)^2$. Either way the suppression is two powers of $\cH/k$, and the limit
$k \gg \cH$ leaves $-k^2\Phi = 4\pi Ga^2\bar\rho\,\delta$.

\emph{(ii)} $(\cH_0/k)^2 = (3.3\times10^{-4}/k)^2$ gives $1.1\times10^{-5}$ at
$k = 0.1\,h\,{\rm Mpc}^{-1}$ and $1.1\times10^{-1}$ at
$k = 10^{-3}\,h\,{\rm Mpc}^{-1}$. So on the scales this course actually uses the
Newtonian field equation is good to a part in $10^5$, while at
$k \sim 10^{-3}$ the neglected term is a ten-percent effect. That is the
low-$k$ edge of the window in the Comment, and this is what closes it.

\emph{(iii)} Because a bound is a claim about the solution, while an identity is
a claim about the equations. The bound $\dot\Phi \lesssim \cH\Phi$ holds for the
growing mode and fails for the decaying one, and one has to check it again in
every new regime --- during a rapid transition, for a modified-gravity model, for
an isocurvature initial condition. The algebraic form has no such exposure: the
time derivative has been eliminated using a second Einstein equation, so the
$\mathcal{O}(\cH/k)$ suppression is visible in the equation itself and holds
wherever the Einstein equations do. Removing an assumption is worth more than
verifying it, because verification does not travel.
\end{solution}

\begin{solution}{Exercise 1.3 (the M\'esz\'aros effect).}
\emph{(i)} With $\delta = 1 + \tfrac32 y$ we have $\delta' = \tfrac32$ and
$\delta'' = 0$, so the left-hand side is
\begin{equation*}
\frac{2+3y}{2y(1+y)}\cdot\frac32 \;-\; \frac{3}{2y(1+y)}\left(1+\frac{3y}{2}\right)
= \frac{1}{y(1+y)}\left[\frac{3(2+3y)}{4} - \frac{3}{2}\left(1+\frac{3y}{2}\right)\right] = 0 ,
\end{equation*}
since both terms in the bracket equal $(6+9y)/4$.

\emph{(ii)} For $y \ll 1$, $\delta \to 1$: constant, no growth at all. For
$y \gg 1$, $\delta \to \tfrac32 y \propto a$, the matter-era growth of
Section~\ref{sec:vlasov}. So the growing mode does nothing during radiation
domination and only switches on at equality. A mode that entered the horizon at
$a_{\rm enter} \ll a_{\rm eq}$ therefore sits idle for the whole interval
between, which is exactly the growth it is missing relative to a mode entering
after equality.

\emph{(iii)} One effect, seen from two sides, and the Poisson equation is what
ties them: $-k^2\Phi = 4\pi Ga^2\bar\rho\,\delta$. During radiation domination
the prefactor $a^2\bar\rho_r$ falls as $a^{-2}$ while $\delta_r$ merely
oscillates, so $\Phi$ decays; and the cold dark matter, with a decaying
potential to fall into and radiation pressure holding the dominant fluid up,
has nothing to collapse toward and freezes. Neither statement causes the other.
They are the same physics written in different variables.

The potential version is the one to turn into a number, because $\Phi \propto
a^{-2}$ is a clean power law and the suppression follows immediately as
$(a_{\rm enter}/a_{\rm eq})^2$. The density version needs the full solution
above, and the amplitude at which $\delta$ freezes depends on how long the
potential took to decay after entry --- which is where the logarithm in
$T \propto k^{-2}\ln k$ comes from, and why it is not visible in the scaling
argument.
\end{solution}

%=====================================================================
\section*{Lecture 2}

\begin{solution}{Exercise 2.1 ($D_2$ in EdS).}
In EdS, $a \propto \tau^2$, $\cH = 2/\tau$, $\Om = 1$, and $D_+ = a$. Try
$D_2 = C a^2 \propto \tau^4$. Then $\dot D_2 = 4C\tau^3$ and
$\ddot D_2 = 12C\tau^2$, while $D_+^2 = a^2 \propto \tau^4$. Normalising
$a = \tau^2$, equation \eqref{eq:D2} reads
\begin{equation*}
12C\tau^2 + \frac{2}{\tau}\,4C\tau^3 - \frac{3}{2}\cdot\frac{4}{\tau^2}\,C\tau^4
= -\frac{3}{2}\cdot\frac{4}{\tau^2}\,\tau^4 ,
\end{equation*}
that is $12C + 8C - 6C = -6$, so $14C = -6$ and $C = -3/7$. Hence
$D_2 = -\tfrac37 D_+^2$, the EdS limit quoted in \eqref{eq:D2}. Note that the
associated growth rate must be defined as $f_2 = \dd\ln|D_2|/\dd\ln a$, cf.\
\eqref{eq:f2_def}: with this sign convention $D_2<0$ and $\ln D_2$ does not
exist. In EdS, $|D_2| \propto a^2$ gives $f_2 = 2$.

The negative sign is a bookkeeping artifact of the convention
$\vPsi^{(2)} = +\nabla_q\varphi^{(2)}$, not a statement that second order opposes
first. Check with a spherical patch, where $\varphi^{(1)}_{,ij} = (\delta_1/3)\delta_{ij}$
and the tidal invariant is $\delta_1^2/3$, so \eqref{eq:2lpt_potential} gives
$\nabla_q\cdot\vPsi^{(2)} = \nabla^2_q\varphi^{(2)} = -\delta_1^2/7$. This has the
\emph{same} sign as $\nabla_q\cdot\vPsi^{(1)} = -\delta_1$: the second-order
displacement reinforces the infall. Equivalently, expanding the spherical
Zel'dovich density $(1-D_+\delta_1/3)^{-3}$ gives $\delta^{(2)} = \tfrac23\delta_1^2$,
whereas the correct second-order result --- the angular average of the
second-order Eulerian kernel $F_2$, derived in the companion notes on Eulerian
perturbation theory --- is $\tfrac{17}{21}\delta_1^2$. Since $17/21 > 2/3$, adding the tidal term makes
collapse \emph{faster} than the straight-line Zel'dovich estimate, not slower.
\end{solution}

\begin{solution}{Exercise 2.2 (the Zel'dovich approximation is exact in 1D).}
In one dimension, Gauss's law says the gravitational field at a point depends
only on the total mass on either side of it. Consider a sheet with Lagrangian
coordinate $q$ displaced to $x = q + \Psi(q,\tau)$. Before shell crossing the
ordering of sheets is preserved, so the mass on either side of our sheet is the
same as it was initially: it does not depend on how far anything has moved.

The peculiar force is therefore fixed by the initial configuration alone. Make
this explicit by integrating the 1D Poisson equation
$\partial_x^2\Phi = \tfrac32\Om\cH^2\delta$ once. Using
$1 + \delta = (1 + \partial_q\Psi)^{-1}$ and $\dd x = (1+\partial_q\Psi)\dd q$,
the mass integral collapses exactly:
\begin{equation*}
\int \delta\,\dd x
= \int \left[\frac{1}{1+\partial_q\Psi} - 1\right](1+\partial_q\Psi)\,\dd q
= -\int \partial_q\Psi\,\dd q = -\Psi ,
\end{equation*}
measuring $\Psi$ from a point where it vanishes. Note that every power of
$\partial_q\Psi$ cancels: this is the algebraic content of the statement that the
enclosed mass does not change. Hence
\begin{equation*}
\partial_x\Phi = -\frac{3}{2}\Om\cH^2\,\Psi(q,\tau) ,
\end{equation*}
with no higher-order corrections. The sign is the physically sensible one: a
sheet displaced to $\Psi < 0$, having fallen inward, feels an inward force
$-\partial_x\Phi < 0$. Substituting into the exact equation of motion
\eqref{eq:newtonian_eom} gives
\begin{equation*}
\ddot\Psi + \cH\dot\Psi - \frac{3}{2}\Om\cH^2\Psi = 0 ,
\end{equation*}
which is exactly the linear growth equation \eqref{eq:growth}. So
$\Psi = D_+(\tau)\,\Psi_0(q)$ solves the full nonlinear problem, and the
Zel'dovich displacement is exact until trajectories first cross.

In three dimensions this fails because the force depends on the full
three-dimensional mass distribution, which the displacement itself changes: the
tidal term of \eqref{eq:2lpt_potential} is precisely the leading correction.
\end{solution}

\end{document}
